\documentclass[11pt,a4paper]{article}
\usepackage[margin=1in]{geometry}
\usepackage{hyperref}
\usepackage{enumitem}
\usepackage{graphicx}
\usepackage{xcolor}

% -----------------------------------------------------
% bvraghav-tiet-ucs503p-article.sty
% -----------------------------------------------------

% Variable: Lab Instructor
\makeatletter \newcommand{\BvrSetLabInstructor}[1]{%
  \gdef\Bvr@LabInstructor{#1}}
% default
\newcommand{\Bvr@LabInstructor}{Dr. Raghav %
  B. Venkataramaiyer}
% public accessor
\newcommand{\BvrLabInstructorValue}{\Bvr@LabInstructor}
\makeatother

% Add subtitle
\usepackage{titling}
% -- Set title bold
\pretitle{\begin{center}\bfseries\fontsize{18bp}{18bp}\selectfont}
\makeatletter
\newcommand{\BvrSubtitle}[1]{%
  \posttitle{%
    \par\end{center}
    \begin{center}\large#1\end{center}
    \vskip 0.5em}%
}
\makeatother

% Hints in the section title(s)
\newcommand{\BvrHint}[1]{%
  {\normalfont\normalsize\itshape\hfill #1}}

% -----------------------------------------------------

\title{MeterIQ: Explainable Smart Meter Analytics Platform}

\BvrSetLabInstructor{Dr. Jeelani}

\BvrSubtitle{%
  Project Proposal for UCS503P  \\
  Submitted to: \BvrLabInstructorValue \\ \bfseries
  Thapar Institute of Engineering and Technology%
}

\author{
  Yashit Arora, CSED%
  \thanks{Roll No: \texttt{1024160006}, %
    Email: \texttt{yarora1\_be24\@thapar.edu}}
  \and
  Aarushi Gahlawat, CSED%
  \thanks{Roll No: \texttt{1024160008}, %
    Email: \texttt{agahlawat\_be24\@thapar.edu}}
}

\date{\today}

\begin{document}
\maketitle

\section{Higher-order goal%
  \BvrHint{\ldots data-driven electricity monitoring}}

This project aims to make smart-meter data easier to analyse and use.
The system will help identify meters whose electricity consumption
looks unusual and will also provide short-term load forecasts.

We will work with smart-meter datasets from Mathura and Bareilly. The
main idea is to bring anomaly detection and load forecasting together
in one platform. A user should be able to search for a meter, look at
its consumption pattern, understand why it may have been flagged, and
see its expected future load.

\section{Time-to-value %
  \BvrHint{\ldots primary heuristic}}

\begin{itemize}[leftmargin=0.7cm]
    \item We will first inspect the Mathura and Bareilly datasets to
    understand the available meter, time and energy-related fields
    before deciding which features and models can be used.

    \item We will then clean the data and prepare the features needed
    for anomaly detection and short-term load forecasting.

    \item Simple baseline approaches will be developed first. This
    will give us a reference point before trying more advanced models.

    \item After the ML part is working, we will connect it with the
    backend, database and dashboard so that the results can be viewed
    through one application.

    \item Git and GitHub will be used throughout the project so that
    changes can be tracked and the work can be developed in smaller
    steps.
\end{itemize}

\section{Problem Statement}

Electricity distribution companies collect large amounts of data from
smart meters, but raw meter readings by themselves are difficult to
use for regular monitoring. A meter may show a sudden increase or
decrease in consumption for many different reasons, and it is not
practical to manually inspect every reading.

The available datasets contain timestamped electricity consumption
along with electrical measurements such as average voltage, average
current and frequency. The data comes from smart meters located in
Mathura and Bareilly and contains readings at a fine time resolution.

The problem addressed by this project is to build a practical system
that can analyse these readings, identify unusual behaviour and
forecast future load. The system should also present the results in a
form that is easier for a user to understand.

A major concern is that an anomaly should not automatically be
labelled as electricity theft. The system therefore focuses on
identifying unusual behaviour and providing supporting information
for further investigation rather than making an unsupported theft
claim.

\section{Proposed Solution %
  \BvrHint{\ldots Software Proposition}}

\subsection{Overview}

We propose a web-based platform called the \textbf{MeterIQ:
Explainable Smart Meter Analytics Platform}. The platform will have
two main parts.

The first part will focus on \textbf{anomaly and theft-risk
detection}. It will study the past consumption behaviour of a meter
and identify patterns that look unusual. The output can be shown as
Normal, Suspicious or High-Risk, along with information that helps
explain the alert.

The second part will focus on \textbf{short-term load forecasting}.
The system will use previous consumption data to predict future load.
The forecast results can then be shown along with the historical
consumption of the meter.

Both parts will be connected to the same application. The user will
not have to work with separate notebooks to see the anomaly and
forecasting results.

\subsection{Core workflow}

\begin{enumerate}[leftmargin=0.7cm]
    \item Smart-meter data from the available datasets is loaded into
    the system.

    \item The data is checked for missing values, duplicates, invalid
    readings, timestamps and sampling issues.

    \item The cleaned data is prepared for analysis and suitable
    features are created based on the actual fields available.

    \item The processed data is given to the anomaly detection and
    load forecasting modules.

    \item The anomaly module produces an anomaly score or flag for
    unusual meter behaviour.

    \item The forecasting module produces future load predictions.

    \item The results are made available through the backend and
    database.

    \item The dashboard allows the user to search for a meter, view
    its consumption history, check its risk level, understand the
    alert and view the forecast.
\end{enumerate}

\subsection{Operational constraints for an educational Software Engineering project}

\begin{itemize}[leftmargin=0.7cm]
    \item The system should remain simple enough to develop and test
    within the project timeline.

    \item The actual Mathura and Bareilly datasets will be inspected
    before deciding which features can be used. We will not assume
    that a feature exists if it is not present in the data.

    \item Forecasting models will be evaluated using chronological
    splits so that future information does not accidentally enter the
    training data.

    \item An anomaly will be treated as an indicator for investigation
    and not as proof of electricity theft.

    \item The frontend, backend, database and ML parts should be kept
    separate enough that they can be developed and tested
    independently.
\end{itemize}

\section{Solution Approach %
  \BvrHint{\ldots Engineering focus}}

The project will combine data analysis, machine learning and normal
software development. We will first understand the data and build
the ML components, and then connect them to the backend, database
and frontend.

\subsection{Data understanding and preprocessing %
  \BvrHint{\ldots data quality}}

The Mathura and Bareilly datasets will first be inspected separately.
The available smart-meter data includes timestamp, energy
consumption, average voltage, average current, frequency and meter
identifier.

The preprocessing stage will include timestamp conversion, checking
the sampling interval, identifying missing values and duplicate
records, checking unusual electrical readings, and validating
meter-wise data coverage.

Time-based features such as hour, day, weekday and month will be
created where useful. Lag and rolling features may also be generated
for forecasting, while ensuring that they use only information that
would have been available at prediction time.

\subsection{Anomaly and risk detection %
  \BvrHint{\ldots ML component}}

The anomaly detection module will focus on finding unusual
consumption behaviour rather than directly classifying a reading as
electricity theft.

The analysis will consider changes in energy consumption together
with contextual variables such as voltage, current, frequency and
time of day. Meter-level behaviour will also be considered because
a pattern that is normal for one meter may be unusual for another.

Candidate unsupervised techniques such as Isolation Forest and Local
Outlier Factor will be evaluated. The final method will be selected
after comparing the stability and usefulness of the detected
anomalies on the available data.

The output will include an anomaly score or flag together with the
timestamp and supporting information that can help a user understand
why the reading was considered unusual.

The system will clearly distinguish between an ML-generated anomaly
and a confirmed theft case. A flagged reading will only indicate
that further investigation may be useful.

\subsection{Short-term load forecasting %
  \BvrHint{\ldots forecasting component}}

The forecasting module will estimate future electricity consumption
from historical meter behaviour.

The project will first establish a simple baseline using previous
historical observations. More advanced machine learning models will
then be compared against this baseline.

Potential features include recent consumption values, lagged
consumption, rolling statistics and calendar information such as
hour, weekday and month. Electrical variables such as voltage,
current and frequency will be considered where they are available
and useful.

Candidate models will include Linear Regression, Random Forest and
gradient-boosting based models. The final model will be selected
based on performance on a chronological test period rather than a
random train-test split.

The main evaluation measures will be MAE, RMSE and MAPE. R-squared
may also be reported where it provides useful information about the
model.

\subsection{Explainability %
  \BvrHint{\ldots understandable results}}

Explainability is an important part of MeterIQ because the purpose
of the system is not only to produce an anomaly flag but also to
make the result understandable.

For each flagged observation, the system will present information
such as the observed consumption, recent historical behaviour,
deviation from normal patterns and relevant electrical measurements.

The interface will distinguish between what the model directly
predicts, what is calculated from the data and what is interpreted
by the user. This avoids presenting a possible explanation as a
confirmed cause when the available data cannot support such a
conclusion.

\subsection{Web architecture %
  \BvrHint{\ldots web-based solution}}

The application will follow a simple full-stack structure.

\begin{itemize}[leftmargin=0.7cm]
    \item \textbf{Frontend:} React.js with HTML and CSS. A charting
    library can be used to display consumption, anomaly and
    forecasting results.

    \item \textbf{Backend API:} Flask or FastAPI will be used to
    provide APIs for meter information, anomaly results, forecasts
    and dashboard data.

    \item \textbf{Data and ML layer:} Python, Pandas, NumPy and
    Scikit-learn will be used for data processing and machine
    learning.

    \item \textbf{Database:} PostgreSQL will be used to store
    application data and the results needed by the dashboard.

    \item \textbf{Version control:} Git and GitHub will be used to
    keep the project organised and track changes.
\end{itemize}

\subsection{CI/CD and fast delivery enabler}

GitHub will be used throughout development. As the project develops,
we plan to add automated testing and build checks so that changes can
be checked before they are merged or deployed.

The CI/CD setup will help us to:

\begin{itemize}[leftmargin=0.7cm]
    \item run tests automatically after changes;
    \item check that the application still builds correctly;
    \item keep deployment steps repeatable; and
    \item make it easier to add new features without breaking the
    existing system.
\end{itemize}

\section{Evaluation Criterion %
  \BvrHint{\ldots measurable and attributable}}

The project will be evaluated using both ML results and the
performance of the software system. We will not claim model
performance before the experiments are actually completed.

\subsection{Primary evaluation metric}

For load forecasting, the main metrics will be:

\begin{itemize}[leftmargin=0.7cm]
    \item \textbf{Mean Absolute Error (MAE):} average absolute
    difference between predicted and actual load.

    \item \textbf{Root Mean Squared Error (RMSE):} gives more weight
    to larger forecasting errors.

    \item \textbf{Mean Absolute Percentage Error (MAPE):} measures
    forecasting error relative to the actual value when suitable
    for the data.

    \item \textbf{R$^2$:} may be used as an additional measure where
    appropriate.
\end{itemize}

For anomaly detection, evaluation will focus on the quality and
usefulness of the detected patterns because confirmed theft labels
are not assumed to be available in the supplied datasets.

\subsection{Secondary evaluation metrics}

The software system will also be checked for:

\begin{itemize}[leftmargin=0.7cm]
    \item correct API responses;
    \item correct retrieval of meter data;
    \item correct display of anomaly and forecast results;
    \item database consistency;
    \item response time where measurable;
    \item integration between the ML and web application layers; and
    \item reliability of the main user workflows.
\end{itemize}

\subsection{Pilot validation plan %
  \BvrHint{\ldots high-level}}

The forecasting models will be evaluated using chronological
validation so that the model is tested on future observations
relative to its training data.

Different forecasting models will be compared using the same
evaluation period. For anomaly detection, the detected patterns
will be inspected at meter level and compared with historical
behaviour to determine whether the results are meaningful and
stable.

Any data leakage discovered during development will be corrected
before the final model comparison.

\section{Scalability %
  \BvrHint{\ldots with foresight}}

The initial analysis will use the supplied Mathura and Bareilly
datasets. The software architecture will be designed so that
additional meters and future readings can be added later without
changing the complete application structure.

The database and API layers will remain separate from the model
code, making it possible to update or retrain models as more data
becomes available.

Future versions may support automated ingestion of new smart-meter
readings and periodic model updates.

\section{Engine availability heuristic}

The availability of a meter for analysis will be assessed using the
completeness and continuity of its recorded observations.

A meter with large gaps in its readings may provide less reliable
input for forecasting and anomaly detection. Therefore, data
availability checks will be performed before using a meter for model
training or detailed investigation.

This measure will be treated as a data-quality indicator and not as
a measure of the physical health of the meter.

\section{Project scope and deliverables %
  \BvrHint{\ldots iteration-friendly}}

\subsection{Initial deliverable %
  \BvrHint{\ldots first iteration}}

The initial deliverable will be a cleaned and documented version of
the supplied Mathura and Bareilly smart-meter data together with
exploratory analysis of consumption and electrical behaviour.

This stage will establish the data quality, time coverage,
meter-wise behaviour and useful features required for the
subsequent modelling stages.

\subsection{Subsequent deliverables}

\begin{itemize}[leftmargin=0.7cm]
    \item A preprocessing and feature-engineering pipeline.
    \item An anomaly detection module for identifying unusual meter
    behaviour.
    \item A short-term load forecasting module.
    \item Evaluation and comparison of candidate models.
    \item An explainability layer for anomaly investigation.
    \item A backend API connecting the analytical modules with the
    application.
    \item A database for storing meter, anomaly and forecast
    information.
    \item A web dashboard for viewing and investigating results.
    \item Documentation and testing of the complete system.
\end{itemize}

\section{Risks and mitigations}

\begin{itemize}[leftmargin=0.7cm]

    \item \textbf{Missing or inconsistent readings:} data-quality
    checks and preprocessing will be performed before modelling.

    \item \textbf{Different behaviour between meters:} the analysis
    will consider meter-level patterns instead of assuming that all
    meters behave in the same way.

    \item \textbf{False anomaly detection:} different approaches
    will be compared and flagged observations will be treated as
    investigation signals rather than confirmed theft.

    \item \textbf{Data leakage:} features will be generated using
    only information available before the prediction time, and
    chronological validation will be used.

    \item \textbf{Model overfitting:} models will be compared on a
    separate chronological evaluation period.

    \item \textbf{Limited ground truth for theft:} the project will
    avoid claiming confirmed theft detection when reliable labels
    are unavailable.

    \item \textbf{System integration:} the frontend, backend,
    database and ML components will be developed as separate
    modules and tested before final integration.
\end{itemize}

\section{Summary}

MeterIQ aims to provide a practical way of working with smart-meter
data from Mathura and Bareilly.

The project starts with understanding and cleaning the available
meter readings and then builds two connected analytical components:
anomaly detection and short-term load forecasting. The anomaly
module is intended to highlight unusual behaviour, while the
forecasting module provides an estimate of expected future
consumption.

The important part of the project is the connection between these
results. Instead of showing an anomaly flag alone, MeterIQ will
allow a user to look at the meter's historical behaviour, understand
the context of the unusual reading and compare it with expected load
behaviour.

The final system will combine the data pipeline, machine learning
models, database, backend APIs and web dashboard into one
application. The project will focus on producing useful and
explainable results while avoiding unsupported claims about
electricity theft.

\end{document}
